\justifying
\NSFCBodyText

\subsubsection{研究背景与意义}

民族地区自然场景图像往往同时包含物体外观和场景文字。路牌、招牌、海报上的维吾尔语既标识地点与商户，也构成图像语义的一部分。现有视觉—语言模型主要在英文或中文图文对上训练，倾向于回答“图中是什么”，对黏着语、从右向左书写、连写变体丰富的维吾尔语场景文字缺乏稳定编码\cite{radford2021learning,chen2023altclip}。若把这类图像只当作一般照片做跨模态检索，文字通道被当成纹理；若只做场景文字识别，又丢掉物体与场景上下文。

本项目面向教育部重点实验室“民族语言多模态提取融合”方向，研究如何把维吾尔语场景文字作为独立视觉语言通道，与图像外观、维吾尔语描述和汉语描述一并编码，并在公开可核验数据上检验跨语言图像—文本检索与小规模图像理解。对象收敛到维—汉图文，不扩展为通用民汉翻译系统、语音交互或舆情治理。

该问题具有方法与应用双重约束。方法上，外观特征、文字区域特征与双语描述的粒度、噪声和脚本方向不同，直接拼接或单一对比损失容易互相干扰。应用上，跨语言检索需要支持汉查询维图、维查询汉图，而不能假设用户与图像文字使用同一语言。项目因此把“第三通道是否有用”写成可消融、可证伪的问题，而不是先承诺平台或语料规模。

实验室开放课题经费有限、周期两年，不适合以大规模预训练或在线系统为交付。可检验的交付是：在公开数据上复现基线、加入文字通道、报告检索与消融。若文字通道只提高 OCR 分数、不提高汉$\leftrightarrow$维检索，则说明识别与理解仍是两件事，项目应收缩主张而不是改写指标。

\subsubsection{国内外研究现状及发展动态分析}

视觉—语言对比学习把图像与句子拉到共享空间，CLIP、ALIGN 与 SigLIP 提供了可迁移的图文编码器\cite{radford2021learning,jia2021scaling,zhai2023sigmoid}。细粒度匹配与统一预训练进一步改善检索和理解，SCAN、ALBEF、BLIP 与 Flamingo 分别强调区域—单词对齐、先对齐后融合、理解—生成一体和少样本条件生成\cite{lee2018stacked,li2021align,li2022blip,alayrac2022flamingo}。早期视觉—语义嵌入则说明跨模态检索的评价应同时报告图像到文本与文本到图像\cite{frome2013devise,faghri2018vse}。这些模型的训练数据以自然描述为主，图像中的文字通常不被单独建模。对维吾尔语场景图，这一隐含假设会带来两类错误：招牌文字被当成纹理噪声，或模型用英文/中文先验去“猜”图意而忽略真实书写。因此，把 CLIP 式编码器直接零样本用于维—汉检索，只能作为对照，不能当作已经解决场景文字问题。

场景文字识别把“图中的字”当作序列转写问题。CRNN 给出端到端卷积—循环基线，ABINet 用双向语言模型迭代校正，TrOCR 与 CLIP4STR 把预训练视觉—语言模型引入识别\cite{shi2017end,fang2021read,li2023trocr,zhao2024clip4str}。评测协议本身会影响结论，因此不能只报单一测试集准确率\cite{baek2019what}。当答案写在图里时，TextVQA、ST-VQA 与 OCR-VQA 表明物体外观不够，必须读取文字\cite{singh2019towards,biten2019scene,mishra2019ocr,hu2020iterative}。OCR-free 文档理解试图跳过显式转写\cite{kim2022ocr}，但其数据仍集中在英文或中文印刷体，不能直接外推到维吾尔语连写招牌。对申请书更关键的是任务接口：STR 输出字符串，图像理解需要的是可与外观特征相加或门控的连续向量。若只把识别结果再翻译成汉语再检索，错误会在 OCR、翻译和检索三级累积，且无法做“去掉文字通道”的干净消融。本项目因此保留显式文字编码器，但把它的输出当作融合特征，而不是停在转写竞赛。

多语图文模型用替换文本编码器、蒸馏或翻译图文对扩展 CLIP，AltCLIP、mCLIP、Chinese CLIP、M3P 与 UC2 在 Multi30k、COCO-CN 等基准上提升跨语言检索\cite{chen2023altclip,chen2023mclip,yang2022chineseclip,ni2021m3p,zhou2021uc2,elliott2016multi30k}。这些评测很少把维吾尔语当作主测试语言。视觉枢纽研究则表明，当句对稀缺时，共享图像可以约束跨语言潜空间\cite{huang2020unsupervised}。针对英—维、汉—维远程语言对，人工翻译的 Multi30k-Distant 显示视觉枢纽对多模态翻译有增益\cite{tayir2024visual,tayir2024unsupervised}。该基准提供了同图维/汉/英描述，但任务是翻译而不是场景文字增强的图像理解。翻译指标对词序和形态变化敏感，维吾尔语作为黏着语更容易被 BLEU 低估或高估。本项目改用检索 Recall@$K$ 与 MRR：它们直接回答“汉查询能否找到对应维语描述的图像”，也便于与无文字通道模型比较。Multi30k 图像本身很少是新疆街景招牌，因此 RUST 仍用于检验文字通道，二者不互相替代。

维吾尔语场景文字已出现可公开使用的数据。SUST 含 60 万合成词级场景文字图，RUST 含 4000 张新疆实景图，后续工作讨论低资源协同编码与视觉预测增强\cite{kong2023sust,yang2024collaborative,ma2023scene}。验收指标仍是识别准确率，文字区域特征尚未进入跨语言检索。MC$^2$ 提供维吾尔语网页单语文本，CUTE 提供中维藏英平行与非平行语料\cite{zhang2024mc2,zhuang2025cute}。CUTE 含机器翻译生成部分，只适合文本预训练，不适合作为图像理解金标准。对开放课题而言，使用这些语料的边界必须写清楚：它们可以改善维语子词切分和文本编码器，但不能用来宣称已经拥有大规模维语图文理解基准。商业平行库不进入本项目数据清单。

申请人前期在知识增强视觉问答、空间共注意力、遥感视觉问答和跨模态检索上形成了多模态表征与对齐基础\cite{yan2024knowledge,yan2022spca,xu2026rshr,xu2026rssr}。本项目把这些能力迁移到维吾尔语场景文字与维—汉图文，而不是重复新能源预测或教学资源构建。

\subsubsection{现有研究不足}

第一，图文检索模型缺少维吾尔语场景文字通道，无法区分“图中是什么”与“图中写了什么”。把多语 CLIP 零样本分数当作维语场景理解能力，会高估资源语言先验、低估文字通道。

第二，维语 STR 工作停在识别准确率，没有检验文字特征是否改善跨语言检索或图像理解。识别正确也不等于融合有用：错误的文字向量可能损害物体检索。

第三，视觉枢纽与 Multi30k-Distant 面向翻译，缺少以检索 Recall@$K$、MRR 和去掉文字通道消融为主的图像理解验收。公开维语 VQA 大规模基准几乎不存在，不能把英文 TextVQA 结论直接写成维语结果。拟建 200--400 题只是辅助协议，主结论仍必须落在公开检索划分上。

\subsubsection{本项目的研究切入点}

项目以“场景文字编码—异构融合—跨语言检索与理解评价”为主线。在 SUST/RUST 上训练并测试维语场景文字编码器，输出供融合使用的区域特征；在外观、文字与维/汉描述之间做分通道编码和门控融合；在 Multi30k-Distant 上报告汉$\leftrightarrow$维互检索，并拟建 200--400 题小规模理解协议。对照包括通用多语 CLIP 零样本、无文字通道对比学习、OCR 后翻译再检索。若去掉文字通道后主指标不下降，或翻译流水线全面不低于视觉枢纽，项目如实报告负结果并收缩主张。
